\section{Pfaffian functions and Pfaffian sets}\label{sec:pfaffian}
This section introduces Pfaffian functions and sets, culminating in Khovanskii's bound on the number of solutions of a system of Pfaffian equations. References for the material in this section are~\cite{khovanskiui1991fewnomials, zell2003quantitative, gabrielov2004complexity}. We assume throughout that $n\in \N_{>0}$.

\subsection{Pfaffian functions}
Pfaffian functions, introduced by Khovanskii~\cite{khovanskiui1991fewnomials}, are functions that satisfy triangular systems of first-order partial differential equations with polynomial coefficients. Pfaffian functions have a well-defined notion of complexity that can be used to prove various quantitative results about the geometric objects they define.

\begin{definition}[Pfaffian function]
Let $\cU \subset \R^n$ be an open set. A \emph{Pfaffian chain} of order $s \in \N$ and chain-degree $\alpha \in \N_{>0}$ over $\cU$ is a sequence of functions $\boldq= (q_1, \dotsc, q_s)$ with $q_i\in C^1(\cU)$ for $i\in [s]$, such that there exist polynomials $P_{ij}\in \R[X_1,\dots,X_n,Y_1,\dots,Y_i]_{\leq \alpha}$, for $i\in [s]$ and $j\in [n]$, that verify
\begin{equation}\label{eq:pfaffchain}
\frac{\partial q_i}{\partial x_j}(x) = P_{ij}(x, q_1(x), \dotsc, q_i(x)).
\end{equation}
A function $g(x)=P(x,q_1(x),\dots,q_s(x))$, with $P\in \R[X_1,\dots,X_n,Y_1,\dots,Y_s]_{\leq \beta}$ for $\beta\in \N$, is called a \emph{Pfaffian function} of chain degree $\alpha$, degree $\beta$, and order $s$.  A function $\cU\to \R^m$ is called Pfaffian if all its components are Pfaffian. 
\end{definition}

The triple $(\alpha, \beta, s)$ is called a \emph{format} of $g$. We denote by $\Pfaff_{\boldq,\beta}(\cU)$ the set of all Pfaffian functions over $\cU$ with chain $\boldq$ and degree $\beta$, and by $\Pfaff_{\alpha,\beta,s}(\cU)$ the set of all Pfaffian functions over $\cU$ with format $(\alpha,\beta,s)$.  When we say that a Pfaffian function $\cU\to \R^m$ has a particular format, we mean that every component has that same format, and we write $\Pfaff_{\alpha,\beta,s}(\cU; \R^m)$.

\begin{remark}
 Note that the Pfaffian chain associated with a Pfaffian function, and hence also a format, is not unique. In particular, if $g\in \Pfaff_{\alpha,\beta,s}(\cU)$, then $g\in \Pfaff_{\alpha',\beta',s'}(\cU)$ for $\alpha'\geq \alpha$, $\beta'\geq \beta$, and $s'\geq s$. In the following, we will often leave the chain implicit, and our results will be stated in terms of the format.
\end{remark}

We call a Pfaffian function $g\colon \cU\to \R$ \emph{autonomous}, if 
\begin{equation}\label{eq:simplepfaff}
  \frac{\partial q_i}{\partial x_j}(x) = P_{ij}(q_1(x),\dots,q_i(x))
\end{equation}
for some $P_{ij}\in \R[Y_1,\dots,Y_i]$, and $g(x)=P(q_1(x), \dots, q_s(x))$ for some $P\in \R[Y_1,\dots,Y_s]$.

\begin{example}\label{ex:1} The following are simple examples of Pfaffian functions.
        \begin{enumerate}
            \item A polynomial $P \in \R[X_1, \dots, X_n]$ gives rise to a Pfaffian function with format $(\alpha,\deg P,0)$ for any $\alpha>0$ via the evaluation homomorphism;
            \item $\exp\colon \R\to \R$ is Pfaffian with format $(1,1,1)$. A Pfaffian chain is $\boldq=(\exp)$ and $P_{11}(X,Y_1)=Y_1$;
            \item $\tanh$ is Pfaffian with format $(2,1,1)$. A Pfaffian chain is $\boldq=(\tanh)$ and $P_{11}(X,Y_1)=1-Y_1^2$;
            \item The logistic sigmoid $\sigma = (1 + \econst^{-x})^{-1}$ is Pfaffian with format $(2, 1, 1)$. A Pfaffian chain is $\boldq=(\sigma)$ and $P_{11}(X,Y_1)=Y_1(1-Y_1)$;
            \item $\arctan$ is Pfaffian with format $(3, 1, 2)$. A Pfaffian chain is $\boldq=((1+x^2)^{-1},\arctan(x))$. The polynomials are $P_{11}(X,Y_1) = -2XY_1^2$ and $P_{21}(X,Y_1,Y_2)=Y_1$. 
        \end{enumerate}
Examples (2-4) are autonomous, while (1) and (5) are not.
\end{example}

%\begin{definition}[Algebraically independent Pfaffian chain]
%\label{defn:algebraically-independent-pfaffian-chain}
%A Pfaffian chain $\boldq = (q_1, \ldots, q_r)$ is called algebraically independent if for all non-zero $P \in \R[Y_1, \ldots, Y_r]$, $P(q_1(x), \ldots, q_r(x))$ is not identically zero.
%\end{definition}

%\begin{remark}
%All the cases in Example~\ref{ex:1} have algebraically independent Pfaffian chains. In~\cite{lotz2024partitioning} it is shown that the chain $\boldq=(1/x,x^m)$ is algebraically independent if $m\in \R\backslash \mathbb{Q}$.    
%\end{remark}

Pfaffian functions enjoy some convenient closure properties under algebraic operations and composition, as illustrated in the following two results (see also~\cite[Proposition 1.8]{zell2003quantitative}).

\begin{lemma}\label{le:pfaffalgebra}
Let $g\in \Pfaff_{\alpha, \beta, s}(\cU)$ and $h\in \Pfaff_{\alpha', \beta', s'}(\cU)$. Then: 
\begin{enumerate}
\item $g + h \in \Pfaff_{\max\{\alpha,\alpha'\}, \max\{\beta,\beta'\},s+s'}(\cU)$;
\item $gh \in \Pfaff_{\max\{\alpha,\alpha'\}, \beta+\beta',s+s'}(\cU)$;
\item For each $i\in [n]$, $\partial g/\partial x_i \in \Pfaff_{\alpha,\alpha+\beta-1,s}(\cU)$. 
\end{enumerate}
\end{lemma}

\begin{proof}
The first two statements are straightforward consequences of the fact that the concatenation of Pfaffian chains is again a Pfaffian chain. For the third statement, assume $g(x)=P(x,q_1(x),\dots,q_s(x))$ for a polynomial $P\in \R[X_1,\dots,X_n,Y_1,\dots,Y_s]_{\leq \beta}$. Then
\begin{equation*}
  \frac{\partial g}{\partial x_i} = \frac{\partial P}{\partial x_i}+\sum_{\ell=1}^s \frac{\partial P}{\partial y_\ell}\frac{\partial q_\ell}{\partial x_i}= \frac{\partial P}{\partial x_i}+\sum_{\ell=1}^s \frac{\partial P}{\partial y_\ell}P_{\ell i}(x,q_1(x),\dots,q_\ell(x)).
\end{equation*}
The resulting expression is a polynomial in $x$ and the $q_j$, and the degree is clearly bounded by $\alpha+\beta-1$.
\end{proof}

\begin{remark}\label{re:1}
Note that the bound on the length of $g+h$ and $gh$ is in general not sharp: if $g$ and $h$ are Pfaffian with respect to the same chain $\boldq$ of length $s$, then clearly the length of $g+h$ and $gh$ is also $s$. It also follows from Lemma~\ref{le:pfaffalgebra} that a linear combination of Pfaffian functions $g_1,\dots,g_m$ with formats $(\alpha_i,\beta_i,s_i)$ is Pfaffian with format $(\max_i \{\alpha_i\},\max_i\{\beta_i\}, \sum_i s_i)$. 
\end{remark}

\begin{lemma}\label{le:composition}
Let $g\in \Pfaff_{\alpha,\beta,s}(\cU;\R^m)$ and $h\in \Pfaff_{\alpha',\beta',s'}(\R^m)$ be Pfaffian functions, and assume $s'\geq 1$. 
\begin{enumerate}
\item $h\circ g \in \Pfaff_{(\alpha'+1)\beta+\alpha-1, \beta\beta', ms+s'}(\cU)$;
\item If $h$ is autonomous, then $h\circ g \in \Pfaff_{\alpha'+\alpha+\beta-1,\beta', ms+s'}(\cU)$;
\item If all the functions in $g$ depend on the same Pfaffian chain $\boldq$ of length $s$, then $h\circ g$ has Pfaffian order $s+s'$.
\end{enumerate}
\end{lemma}

\begin{proof}
Let $g=(g_1,\dots,g_m)$, and for each $i\in [m]$, let $\boldq^{(i)}=(q_1^{(i)},\dots,q_s^{(i)})$ be a Pfaffian chain for $g_i$. Let $\boldq'=(q_1',\dots,q_{s'}')$ be a Pfaffian chain for $h$. Let $P$ and $P_i$, $i\in [m]$, be the polynomials representing $h$ and the $g_i$, respectively. The composition $h\circ g$ is then represented as 
\begin{equation}\label{eq:longpoly}
 (h\circ g)(x) = P(P_1(x,\boldq^{(1)}(x)),\dots,P_m(x,\boldq^{(m)}(x)),\boldq'(g(x))),
\end{equation}
which is clearly a polynomial in $x=(x_1,\dots,x_n)$ and in the chain
\begin{equation}\label{eq:newchain}
  (q_1^{(1)},\dots,q_s^{(1)},\dots,q_1^{(m)},\dots,q_s^{(m)},q'_1\circ g,\dots,q_{s'}'\circ g)
\end{equation}
of length $ms+s'$. It remains to be seen that this is a Pfaffian chain. For the first $ms$ functions in this chain there is nothing to show. It remains to be seen that the derivatives of $q_i'\circ g$ can be expressed as polynomials in $x$ and in the previous elements of the chain. In fact, for $j\in [n]$ we have
\begin{equation}\label{eq:composite}
  \frac{\partial (q'_i\circ g)}{\partial x_j} = \sum_{\ell=1}^m \frac{\partial q'_i}{\partial x_\ell}\frac{\partial g_\ell}{\partial x_j} = \sum_{\ell=1}^m P'_{i\ell}(g(x),q_1'\circ g(x),\dots,q_i'\circ g(x))  \frac{\partial g_\ell}{\partial x_j},
\end{equation}
where the $P_{i\ell}'$ are the polynomials associated to the chain $\boldq'$. 
This shows that the partial derivative is a polynomial in $q_1',\dots,q_i'$. 
Note that $g_{\ell}$ and, by virtue of Lemma~\ref{le:pfaffalgebra}(3), its partial derivatives are polynomials in $q_1^{(\ell)},\dots,q_s^{(\ell)}$. This establishes that~\eqref{eq:newchain} is a Pfaffian chain of length $ms+s'$. 

To prove the degree bounds in the general case (1), note that the degree of the composition of polynomials is bounded by the product of the degrees. The degree of $P_{i\ell}'(g(x),q_1',\dots,q_i')$ is bounded by $\alpha'\beta$, and Lemma~\ref{le:pfaffalgebra}(3) shows that the degree of $\partial g_{\ell}/\partial x_i$ is bounded by $\alpha+\beta-1$, establishing the bound $(\alpha'+1)\beta+\alpha-1$ on the degree of the new Pfaffian chain~\eqref{eq:newchain}. 
Moreover, the degree of the polynomial~\eqref{eq:longpoly} representing $h\circ g$ is clearly bounded by $\beta\beta'$. 
For the autonomous case (2), note that the degree of~\eqref{eq:composite} is $\alpha'+\alpha+\beta-1$ if the $P_{ij}'$ only depend on the elements $q_1',\dots,q_i'$ and not explicitly on $x$, since $P'_{ij}(q_1'\circ g,\dots,q_i'\circ g)$ has degree $\alpha'$ in the last $s'$ chain elements while $\partial g_\ell/\partial x_j$ has degree $\alpha+\beta-1$ in the remaining variables. We also observe from~\eqref{eq:longpoly} that in the autonomous case, $h\circ g$ depends only on the last $s'$ elements of its Pfaffian chain in the same way that $h$ depends on its Pfaffian chain $\boldq'$. This implies that the degree is bounded by $\beta'$.
Case (3) is clear.
\end{proof}

\begin{remark}\label{re:affine}
The special case where $g\colon \R^n\to \R^m$ is a linear map shows that the Pfaffian structure is invariant under affine change of coordinates.
\end{remark}

We are also interested in the geometric objects defined by Pfaffian functions.

\begin{definition}
A \emph{Pfaffian set} (or Pfaffian variety) is a set of the form 
\begin{equation*}
  V = \mathcal{Z}(g_1, \dotsc, g_k) = \{x\in \R^n\colon g_1(x)=\cdots = g_k(x)=0\},
\end{equation*}  
where the $g_i\colon \R^n\to \R$ are Pfaffian functions. A {\em basic semi-Pfaffian set} is a set of the form
\begin{equation*}
 B = \{x\in \R^n \colon g_1(x)=\cdots = g_k(x)=0, h_1(x)>0,\dots,h_{\ell}(x)>0\},
\end{equation*}
where the $g_i$ and $h_j$ are Pfaffian functions. A {\em semi-Pfaffian set} is a finite union of basic semi-Pfaffian sets.
\end{definition}

Note that semi-Pfaffian sets are precisely the sets that can be written as unions, intersections, and complements of sets defined by expressions of the form $g_i(x) \star 0$, with $\star \in \{=,\, \leq,\, \geq\}$. 

\begin{remark}
  Since the concatenation of Pfaffian chains is again a Pfaffian chain, we can assume without loss of generality that the functions $g_i$ and $h_j$ in the definition of a (semi-)Pfaffian set are Pfaffian with respect to the same Pfaffian chain.
\end{remark} 

We conclude this subsection with the following central result on Pfaffian systems.
We call a solution $a\in \R^n$ of a system of equations $F(x)=0$ regular (or non-degenerate), if the differential $\diff{F}(a)$ has maximal rank. 

\begin{theorem}[Khovanskii] \label{thm:khovanskii}
Let $F=(f_1,\dots,f_n) \colon \R^n\to \R^n$ be a Pfaffian function with chain $\boldq=(q_1,\dots,q_s)$ and component-wise formats $(\alpha, \beta_i, s)$, for which the functions in the Pfaffian chain depend only on $\{x_1,\dots,x_k\}$ for some $k\leq n$. Then the number of regular real solutions of the system $F(x) = 0$ is bounded by 
\begin{equation*}
2^{\frac{s(s - 1)}{2}} \beta_1 \dotsm \beta_n \left(\beta_1 + \dotsb + \beta_n - n + \min\{k, s\}\alpha + 1 \right)^s.
\end{equation*}
\end{theorem}

In~\cite[\S 3.12, Corollary 5]{khovanskiui1991fewnomials}, a version of Theorem~\ref{thm:khovanskii} is derived from more general results. The standard formulations (e.g.,~\cite[Theorem 3.1]{gabrielov2004complexity}) use $\min\{n,s\}$ in place of $\min\{k,s\}$; the refined bound follows from the proof in~\cite{khovanskiui1991fewnomials}, where the Rolle-type induction only involves the $k$ variables on which the chain depends. The following result is a slightly sharper version of~\cite[Corollary 3.3]{gabrielov2004complexity} and~\cite[Theorem 2.3]{jones2012density}, using the refined bound of Theorem~\ref{thm:khovanskii}.

\begin{corollary}
Let $F=(f_1,\dots,f_m) \colon \R^n \to \R^m$ be a Pfaffian function with chain $\boldq=(q_1,\dots,q_s)$ and formats $(\alpha, \beta, s)$, for which the functions in the Pfaffian chain depend only on $k\leq n$ of the variables. Then the number of connected components of the zero set $\mathcal{Z}(f_1,\dots,f_m)$ is bounded by 
\begin{equation}\label{eq:jones}
2^{\frac{s(s - 1)}{2}+1} (\alpha+2\beta-1)^{n-1}\beta \left( n (\alpha+2\beta-2)+\alpha(\min\{k, s\}-1) + 2 \right)^s.
\end{equation}
\end{corollary}

The proof follows along the lines of the proof of Proposition 11.5.2 and 11.5.3 in~\cite{bochnak2013real}, using Lemma~\ref{le:pfaffalgebra}(3) to bound the degrees of the derivatives and invoking 
Theorem~\ref{thm:khovanskii} instead of B\'ezout's theorem. While the bound~\eqref{eq:jones} is not as sharp as possible, it gives a convenient way of counting the number of solutions to a system of Pfaffian
equations without reference to the number of equations.

